\documentclass[11pt]{article}
\usepackage[a4paper,margin=1in]{geometry}
\usepackage[T1]{fontenc}
\usepackage{lmodern}
\usepackage{amsmath,amssymb,amsthm,mathtools}
\usepackage[hidelinks]{hyperref}
\usepackage{parskip}

\theoremstyle{plain}
\newtheorem{theorem}{Theorem}[section]
\newtheorem{lemma}[theorem]{Lemma}
\newtheorem{proposition}[theorem]{Proposition}
\newtheorem{corollary}[theorem]{Corollary}
\theoremstyle{definition}
\newtheorem{definition}[theorem]{Definition}
\newtheorem{assumption}[theorem]{Assumption}
\newtheorem{example}[theorem]{Example}
\theoremstyle{remark}
\newtheorem*{remark}{Remark}

\newcommand{\R}{\mathbb{R}}
\newcommand{\norm}[1]{\left\lVert#1\right\rVert}
\newcommand{\abs}[1]{\left\lvert#1\right\rvert}
\newcommand{\trc}{\mathsf{T}}
\DeclareMathOperator*{\argmin}{arg\,min}
\DeclareMathOperator{\rank}{rank}
\DeclareMathOperator{\spn}{span}

\title{\bfseries Ensuring and Certifying the Global\\
Empirical-Acceleration Estimate\\[4pt]
\large A Self-Contained Mathematical Treatment}
\author{}
\date{\today}

\begin{document}
\maketitle

\begin{abstract}
We study when the empirical acceleration (EA) recovered from noisy
observations of a forced trajectory is \emph{the global} minimizer of the
estimation objective, and --- the central question --- when this can be
\emph{proved}. The problem is a regularized nonlinear least-squares
inverse problem constrained by an ordinary differential equation. We
show that an unconditional, a-priori global certificate does not exist for
the operational problem (a consequence of computational complexity,
dimension, and identifiability, not of any single proof technique), but
that a \emph{conditional, constructive, per-run} certificate does. The
treatment is self-contained and develops every notion --- including the
tangential cone condition --- from first principles, with proofs.
Multisensor fusion is treated in detail and shown to be the setting in
which the certificate is most attainable.
\end{abstract}

\section*{How to read this document}
The audience is a reader with a mathematics and numerical-methods
background but no prior exposure to inverse-problem regularity conditions.
\textbf{Proof policy.} Results we prove in full are stated as
Theorem/Lemma/Proposition with a \emph{Proof}. Results that are standard
but whose complete proofs are long are stated precisely with a
\emph{Proof sketch} and a citation to a full treatment; no result is
asserted without either a proof or a citation. The operative certificates
(Theorems~\ref{thm:convex}, \ref{thm:path}, the impossibility
Theorem~\ref{thm:cr}, and the multisensor Lemma~\ref{lem:stack}) are
proved in full.

\section{The problem and its operator}\label{sec:prob}

Let the state $x(t)=(r(t),v(t))\in\R^{6}$ evolve by
\begin{equation}
\dot r = v,\qquad \dot v = g(r,t) + a(t),
\label{eq:ode}
\end{equation}
where $g:\R^{3}\times[0,T]\to\R^{3}$ is a known, real-analytic force
field (gravity, third bodies, radiation), and $a:[0,T]\to\R^{3}$ is the
unknown empirical acceleration. Write the flow
$x(t)=\Phi_{0\to t}(x_0;a)$: the solution operator of \eqref{eq:ode},
returning the state at time $t$ obtained by integrating from the
initial condition $x_0=(r_0,v_0)$ under the driving input $a$. The
semicolon separates the initial condition from the input function on
which the flow depends, and the subscript $0\to t$ records the
integration interval; thus $\Phi_{0\to t}(\,\cdot\,;a):\R^{6}\to\R^{6}$
for each fixed $a$ and each fixed $t\in[0,T]$. We emphasize the
structural fact used throughout:

\begin{remark}[Affine input]\label{rem:affine}
$a$ enters \eqref{eq:ode} \emph{linearly}. All nonlinearity of
$a\mapsto\Phi$ originates in $g$'s dependence on the state, propagated by
the flow; none originates in $a$ itself.
\end{remark}

Expand $a$ in a fixed degree-$d$ B-spline basis on knot spacing $h$,
$a(t)=\sum_{k=1}^{n}c_k B_k(t)$, $c\in\R^{p}$, $p=3n$. With $a$ now
parametrized by the coefficient vector $c$, we write
$\Phi_{0\to t}(x_0;c)$ for $\Phi_{0\to t}(x_0;a[c])$. Observations are
$y_i=\mathcal H\big(\Phi_{0\to t_i}(x_0;c)\big)+\varepsilon_i$,
$i=1,\dots,N$, $\varepsilon_i\sim\mathcal N(0,\Sigma_i)$, with sensor
map $\mathcal H:\R^{6}\to\R^{k}$ converting a full state into the
quantity an individual sensor reports --- concretely
$(\rho,\mathrm{az},\mathrm{el},\dot\rho)$ for a radar
($k=4$), $(\mathrm{az},\mathrm{el})$ for an optics measurement
($k=2$), the position component $r$ for a position-fix ($k=3$), or
the specific force $g(r,t)+a(t)$ for an accelerometer ($k=3$). Thus
$\mathcal H\big(\Phi_{0\to t_i}(x_0;c)\big)$ is the model-predicted
observation at $t_i$ given EA coefficients $c$. Stacking and
pre-whitening, define the forward operator $F:\R^{p}\to\R^{m}$,
\begin{equation}
F(c)_i=\Sigma_i^{-1/2}\big(y_i-\mathcal H(\Phi_{0\to t_i}(x_0;c))\big).
\label{eq:F}
\end{equation}
$F$ is $C^{2}$ on any compact set on which the trajectory avoids the
singularities of $g$ (here $r\neq0$); we make this precise in
\S\ref{sec:local}. The inverse problem ``recover $a$'' is
ill-posed in the sense of Hadamard \cite{tikhonov1977,engl1996}: small
data perturbations need not depend continuously on $c$ without
regularization.

\section{The optimization and the meaning of ``global''}\label{sec:opt}

Let $\Omega=L^{\trc}L\succeq0$ be the order-$m$ spline roughness matrix.
The estimator solves
\begin{equation}
\hat c(\lambda)\in\argmin_{c\in\mathcal C}\;
\Psi_\lambda(c),\qquad
\Psi_\lambda(c)=\tfrac12\norm{F(c)}^{2}+\tfrac{\lambda}{2}\,c^{\trc}\Omega c,
\label{eq:Psi}
\end{equation}
over the compact feasible set $\mathcal C=\{\,c:\norm{a[c]}_\infty\le
A\,\}$, a box of known diameter $D$ (boundedness of EA plus basis
conditioning). ``The global EA'' means a minimizer of $\Psi_{\lambda}$ at
the operational $\lambda$ (or, in the limit, of $\Psi_0$ restricted to
$\mathcal C$). \emph{Whether this minimizer is unique} is itself a
theorem, not an assumption (\S\ref{sec:ident}).

\begin{remark}
Magnitude bounds on $g$ (truncated spherical-harmonic series with bounded
coefficients), on third-body parameters, and on the state
($r\in[r_{\min},r_{\max}]$, $r_{\min}>0$) bound the \emph{curvature} of
$F$ (\S\ref{sec:local}). They do not bound the \emph{number or
arrangement} of stationary points of $\Psi_\lambda$; that is a separate,
landscape question, which is the source of all difficulty below.
\end{remark}

\section{Nonconvexity: a worked scalar model}\label{sec:toy}

Before the general theory, a one-line model exhibits the entire
mechanism. Let $F(c)=\sin(\omega c)-y$, $c,y\in\R$, mimicking a forward
map whose nonlinearity scale is $\omega$ (the analogue of the flow's
sensitivity amplification, $\omega\leftrightarrow L_gT$ of
\S\ref{sec:local}).

\begin{example}\label{ex:scalar}
Stationarity of $\tfrac12(\sin\omega c-y)^{2}$ requires $\sin\omega c=y$
or $\cos\omega c=0$. On an interval of length $\ell$ each equation has
$\Theta(\omega\ell/\pi)$ roots, so the misfit has only $O(1)$ stationary
points when $\omega\ell\lesssim\pi$ but $\Theta(\omega\ell)$ of them once
$\omega\ell\gg\pi$. Thus \emph{boundedness of the data ($\abs y\le1$) and
of every derivative does not bound the number of minima}; only the
product $\omega\ell$ --- a nonlinearity-times-domain quantity --- does.
That same product is precisely $L_gT$ in the orbital problem
(\S\ref{sec:local}): the take-away the reader should carry forward is that
controlling forces controls $\omega$, never $\ell$, and certifiability
lives in their product.
\end{example}

This is why ``bound the forces'' alone cannot certify the global
minimum: it controls $\omega$ (curvature), not $\ell$ (the set we must
search) nor the count of basins their product produces.

\section{Local theory: Gauss--Newton and the convergence ball}
\label{sec:local}

Let $J(c)=F'(c)\in\R^{m\times p}$, the Jacobian of the pre-whitened
residual operator with respect to $c$. Differentiating
\eqref{eq:F}, block-row $i$ factors by the chain rule as
\begin{equation}
J(c)_i \;=\; -\,\Sigma_i^{-1/2}\,
\underbrace{\frac{\partial \mathcal H}{\partial x}\big(\Phi_{0\to t_i}(x_0;c)\big)}_{\text{sensor linearization}\;\in\R^{k\times 6}}\;
\underbrace{\frac{\partial \Phi_{0\to t_i}}{\partial c}(x_0;c)}_{\text{state-transition sensitivities}\;\in\R^{6\times p}},
\label{eq:Jfactor}
\end{equation}
the second factor obtained by integrating the variational
(state-transition) equation along the trajectory --- hence the
qualifier ``variational/state-transition Jacobian''. Let
$G(c)=J(c)^{\trc}J(c)\succeq0$. Gauss--Newton, regularized
(IRGN), iterates
\begin{equation}
c_{k+1}=c_k-\big(G_k+\lambda_k\Omega\big)^{-1}
\big(J_k^{\trc}F(c_k)+\lambda_k\Omega c_k\big),\quad
\lambda_k=\max(\lambda_{\min},\lambda_0 q^{k}).
\label{eq:irgn}
\end{equation}

\begin{lemma}[Bounded input gives explicit curvature]\label{lem:KT}
Assume \textup{(A1)} $\norm{a}_\infty\le A$ on $[0,T]$ and \textup{(A2)}
$r(t)\in[r_{\min},r_{\max}]$, $r_{\min}>0$, on $[0,T]$. Then $F\in
C^{2}(\mathcal C)$ and there is a finite, in-principle computable constant
\begin{equation}
\sup_{\mathcal C}\norm{F''}\;\le\;K_T\;=\;C\,e^{\,\beta L_gT},\qquad
L_g=\sup_{r\in[r_{\min},r_{\max}],\,t}\norm{\partial g/\partial r},
\label{eq:KT}
\end{equation}
for a constant $\beta\le2$ and $C$ depending on $\norm B$, $T$ and the
basis conditioning. (The constant in the exponent is $\le2$ because the
\emph{second} variation of the flow carries the state-transition matrix
quadratically; only the qualitative fact ``exponential in $L_gT$''
matters below.) Moreover, if the unforced trajectory stays a fixed
margin above the singular set on $[0,T]$ and $A$ is small enough that the
Gr\"onwall deviation does not consume that margin, then \textup{(A1)}
yields a compact tube and hence \textup{(A2)} with $r_{\min}>0$ --- the
situation on any arc cut short of reentry.
\end{lemma}

\begin{proof}[Sketch]
By Remark~\ref{rem:affine} the second derivative of $c\mapsto\Phi$ has no
contribution from the input; differentiating the variational equation
twice, $F''$ is a time-integral of $\partial^{2}g/\partial r^{2}$
contracted with the state-transition matrix. On $[r_{\min},r_{\max}]$,
$\partial^{2}g/\partial r^{2}$ (essentially $\partial^{3}V/\partial
r^{3}$ plus bounded harmonics) is bounded; the state-transition matrix
satisfies $\norm{M(t)}\le e^{L_gT}$ by Gr\"onwall's inequality applied to
the linear variational system \cite{hairer1993,ortega1970}. The second
variation pairs \emph{two} transition matrices with one bounded
$\partial^{2}g/\partial r^{2}$, hence is controlled by $e^{\beta L_gT}$,
$\beta\le2$; composing with the bounded $\mathcal H$ gives \eqref{eq:KT}.
The compact-tube claim is a second Gr\"onwall estimate on the deviation
from the unforced trajectory, valid while that deviation does not reach
the singular set.
\end{proof}

\begin{proposition}[Convergence ball; regularization enlarges it]
\label{prop:ball}
Under the affine-covariant Lipschitz condition
$\norm{(G+\lambda\Omega)^{-1}(J(c')-J(c))(c'-c)}\le\omega(\lambda)
\norm{c'-c}^{2}$, the iteration \eqref{eq:irgn} converges q-linearly to a
stationary point of $\Psi_\lambda$ from any $c_0$ with
$\norm{c_0-\hat c}<2/\omega(\lambda)$ \cite{deuflhard2004,nocedal2006}.
Furthermore
\begin{equation}
\rho(\lambda):=\frac{2}{\omega(\lambda)}\;\ge\;
\frac{2\lambda\,\sigma_{\min}(\Omega)}{K_T}>0 .
\label{eq:rho}
\end{equation}
\end{proposition}

\begin{proof}
The first statement is the affine-covariant Newton--Kantorovich theorem
\cite{deuflhard2004} for the regularized normal map. For
\eqref{eq:rho}: $G\succeq0$ and Weyl's inequality give
$\sigma_{\min}(G+\lambda\Omega)\ge\lambda\,\sigma_{\min}(\Omega)$, hence
$\norm{(G+\lambda\Omega)^{-1}}\le(\lambda\sigma_{\min}(\Omega))^{-1}$.
With $\norm{J(c')-J(c)}\le K_T\norm{c'-c}$
(Lemma~\ref{lem:KT}) the affine-covariant constant obeys
$\omega(\lambda)\le K_T/(\lambda\sigma_{\min}(\Omega))$; invert.
\end{proof}

A word on what Proposition~\ref{prop:ball} does and does not say, since it
is easy to over-read. It is the \emph{local} affine-covariant
Gauss--Newton convergence ball: it presumes the iteration is already near
\emph{a} stationary point with a residual small relative to the curvature
(the standard small-residual Deuflhard hypothesis \cite{deuflhard2004});
it does not assert convergence from arbitrary starts, nor that the
stationary point reached is the global one. Equation \eqref{eq:rho} is
rigorous \emph{with an explicit constant} (not merely in form), but
$K_T\sim e^{\beta L_gT}$ is exponentially conservative
(\S\ref{sec:toy}'s $\omega\ell$). The useful content is the
\emph{mechanism}: more regularization provably enlarges the ball. Whether
the ball is operationally large, and whether the landscape is unimodal,
are separate questions taken up in
\S\S\ref{sec:tcc}--\ref{sec:ceiling}.

\section{The tangential cone condition}\label{sec:tcc}

This section is the conceptual center of the document, so we proceed
slowly. We have a basin around \emph{a} stationary point
(\S\ref{sec:local}); what we actually want is the assurance that the
stationary point we reach is \emph{the global one} --- equivalently, that
the misfit has \emph{no spurious traps} on the region we search. The
tangential cone condition (TCC) is exactly the structural hypothesis that
delivers this, and it is the inverse-problems analogue of the
``every critical point is global'' conditions familiar from convex and
gradient-dominated optimization. A reader meeting it for the first time
should focus on the geometric reading given after the definition.

\begin{definition}[Tangential cone condition, TCC]\label{def:tcc}
$F$ satisfies the TCC with constant $\eta\in(0,1)$ on a ball
$B_\rho(c^{\dagger})$ if for all $c,\tilde c\in B_\rho(c^{\dagger})$,
\begin{equation}
\norm{F(c)-F(\tilde c)-F'(c)(c-\tilde c)}\;\le\;
\eta\,\norm{F(c)-F(\tilde c)} .
\label{eq:tcc}
\end{equation}
\end{definition}

\paragraph{Why this and not ordinary smoothness.} A $C^{2}$ map obeys
$\norm{F(c)-F(\tilde c)-F'(c)(c-\tilde c)}\le\tfrac12 K_T
\norm{c-\tilde c}^{2}$. For ill-posed $F$ the right-hand side can dwarf
$\norm{F(c)-F(\tilde c)}$ (small singular values make $F$ move little
while $c$ moves a lot), so quadratic smoothness gives no control of the
\emph{misfit landscape}. TCC instead measures the linearization error
\emph{relative to how much $F$ actually moves}: it forbids the image of a
segment from leaving a cone about its tangent --- geometrically, the
surface never folds back on itself.

\begin{lemma}[TCC removes spurious critical points]\label{lem:tcc}
If $F$ satisfies the TCC with $\eta<1$ on $B_\rho(c^\dagger)$ and
$F(c^\dagger)=0$ (attainable, noise absorbed), then every stationary
point of $c\mapsto\tfrac12\norm{F(c)}^{2}$ in $B_\rho(c^\dagger)$ is a
global zero; in particular the minimizer is unique up to
$\ker F'(c^\dagger)$.
\end{lemma}

\begin{proof}
The idea is to project away the linear (tangent) term in the TCC and let
stationarity kill the remaining component. Let $c^\star\in B_\rho$ be
stationary, $F'(c^\star)^{\trc}F(c^\star)=0$, and let $P$ be the
orthogonal projector onto $\operatorname{range}F'(c^\star)$. Stationarity
says precisely $F(c^\star)\in\big(\operatorname{range}F'(c^\star)\big)^{
\perp}$ (because $F'^{\trc}F=0\iff F\perp\operatorname{range}F'$), hence
\[
P\,F(c^\star)=0,\qquad\text{so}\qquad (I-P)\,F(c^\star)=F(c^\star).
\]
Apply the TCC \eqref{eq:tcc} with $c=c^\star$, $\tilde c=c^\dagger$ and
$F(c^\dagger)=0$:
\[
\norm{F(c^\star)-F'(c^\star)(c^\star-c^\dagger)}\le
\eta\,\norm{F(c^\star)} .
\]
Now $F'(c^\star)(c^\star-c^\dagger)\in\operatorname{range}F'(c^\star)$,
so it is annihilated by $I-P$. Since $I-P$ is an orthogonal projector
($\norm{I-P}\le1$),
\[
\norm{F(c^\star)}
=\norm{(I-P)F(c^\star)}
=\big\|(I-P)\big(F(c^\star)-F'(c^\star)(c^\star-c^\dagger)\big)\big\|
\le\eta\,\norm{F(c^\star)} .
\]
With $\eta<1$ this forces $\norm{F(c^\star)}=0$: every stationary point
in the ball is a global zero. Uniqueness modulo $\ker F'$: if
$F(c_1)=F(c_2)=0$ with $c_1,c_2\in B_\rho$, the TCC with $c=c_1$,
$\tilde c=c_2$ gives $\norm{F'(c_1)(c_1-c_2)}\le\eta\cdot0=0$, so
$c_1-c_2\in\ker F'(c_1)$.
\end{proof}

\begin{remark}[The noisy case is the honest one]
The hypothesis $F(c^\dagger)=0$ is the \emph{attainable} idealization
(the model can fit the data exactly). With genuine noise no exact zero
exists; repeating the projection argument with $F(c^\dagger)\neq0$ yields
$\norm{F(c^\star)}\le(1-\eta)^{-1}\norm{F(c^\dagger)}$, i.e.\ spurious
critical points sit within an $O(\text{noise level})$ band of the
optimum. This is exactly the ``semiconvergence'' picture: one stops the
iteration at the noise level (discrepancy principle), and TCC guarantees
no \emph{worse-than-noise} trap in the ball \cite{kaltenbacher2008}.
\end{remark}

\begin{theorem}[IRGN under TCC]\label{thm:irgn}
If $F$ satisfies the TCC with $\eta<1$ on a ball containing $c_0$ and
$c^\dagger$, then \eqref{eq:irgn} (and Landweber, Levenberg--Marquardt)
converges to a solution; with a source condition the convergence is
order-optimal.
\end{theorem}
\begin{proof}[Sketch]
This is the central result of nonlinear-ill-posed regularization theory;
see Scherzer \cite{scherzer1995}, Hanke--Neubauer--Scherzer
\cite{hanke1995}, and the monograph treatment of IRGN by
Kaltenbacher--Neubauer--Scherzer \cite{kaltenbacher2008}
(Bakushinskii's method \cite{bakushinskii1992}). The mechanism is
Lemma~\ref{lem:tcc}: TCC makes $\norm{F}$ a global merit function on the
ball, so the regularized Gauss--Newton step is a descent direction
without spurious traps.
\end{proof}

\begin{remark}[Relation to Polyak--{\L}ojasiewicz]
TCC is the inverse-problems form of the gradient-domination /
Polyak--{\L}ojasiewicz inequality
$\norm{\nabla\Psi_0}^{2}\ge2\mu(\Psi_0-\Psi_0^\star)$
\cite{polyak1963,lojasiewicz1963,karimi2016}: both express ``every
critical point is global''. The TCC implies a (local) gradient-domination
estimate: writing $\Delta=c-c^\dagger$ and using \eqref{eq:tcc} with
$F(c^\dagger)=0$ gives $\langle F(c),F'(c)\Delta\rangle\ge(1-\eta)
\norm{F(c)}^{2}$, hence
$\norm{\nabla\Psi_0(c)}\,\norm{\Delta}\ge\langle F'(c)^{\trc}F(c),
\Delta\rangle\ge 2(1-\eta)\big(\Psi_0(c)-\Psi_0^\star\big)$,
a P{\L}-type inequality on the ball (the precise modulus depends on
$\norm\Delta$ and is not a universal constant; see
\cite{karimi2016} for the general correspondence). Neither condition is
implied by magnitude bounds; both fail in the same regimes.
\end{remark}

\section{Certificate I: a checkable convex regime}\label{sec:convex}

TCC is rarely verifiable a priori. The following is a fully constructive
substitute for the \emph{regularized} problem.

\begin{theorem}[Global strong convexity threshold]\label{thm:convex}
Assume \textup{(A1)}--\textup{(A2)} and let $R=\sup_{\mathcal
C}\norm{F}<\infty$ (finite by Lemma~\ref{lem:KT}). If
\begin{equation}
\lambda\;>\;\lambda_{\mathrm{cvx}}\;:=\;\frac{R\,K_T}{\sigma_{\min}(\Omega)},
\label{eq:lamcvx}
\end{equation}
then $\Psi_\lambda$ is strongly convex on the convex set $\mathcal C$
with modulus $\lambda\sigma_{\min}(\Omega)-RK_T>0$. Consequently
$\Psi_\lambda$ has a \emph{unique} minimizer on $\mathcal C$, and any
descent method (gradient descent, Gauss--Newton, Levenberg--Marquardt)
from \emph{any} initial point converges to it q-linearly.
\end{theorem}

\begin{proof}
$\nabla^{2}\Psi_\lambda(c)=G(c)+\sum_{i}F_i(c)\nabla^{2}F_i(c)
+\lambda\Omega$. The middle term has norm
$\le\norm{F(c)}\sup_i\norm{\nabla^{2}F_i}\le R\,K_T$ on $\mathcal C$
(Lemma~\ref{lem:KT}). Since $G\succeq0$ and, by Weyl,
$\lambda_{\min}(\lambda\Omega)=\lambda\sigma_{\min}(\Omega)$,
\[
\nabla^{2}\Psi_\lambda \succeq
\big(\lambda\sigma_{\min}(\Omega)-RK_T\big)\,I \succ 0
\]
whenever \eqref{eq:lamcvx} holds. Strong convexity on the convex
$\mathcal C$ gives uniqueness; global q-linear convergence of descent
methods from any start is the standard strongly-convex guarantee
\cite[Ch.~2]{nesterov2004}.
\end{proof}

\begin{corollary}[A surrogate global certificate]\label{cor:surr}
For $\lambda\ge\lambda_{\mathrm{cvx}}$ the global minimizer of
$\Psi_\lambda$ is computable with a \emph{single} solve and a
\emph{certificate} (positive definiteness of $\nabla^{2}\Psi_\lambda$,
checkable by one smallest-eigenvalue estimate). No landscape assumption
and no TCC is needed.
\end{corollary}

The price: this is the global of the \emph{heavily regularized} surrogate,
biased toward the prior. The operational object is
$\hat c(\lambda_{\min})$.

\section{Certificate II: continuation and per-run certification}
\label{sec:path}

Define the solution path $\lambda\mapsto c^\star(\lambda)$ by the
stationarity map
$\Theta(c,\lambda):=J(c)^{\trc}F(c)+\lambda\Omega c=0$.

\begin{theorem}[Per-run path certificate]\label{thm:path}
Suppose $c^\star(\lambda_{\mathrm{cvx}})$ is the unique minimizer of
Theorem~\ref{thm:convex}, and along a continuation
$\lambda:\lambda_{\mathrm{cvx}}\to\lambda_{\min}$ the regularized Hessian
satisfies $\sigma_{\min}\big(\nabla^{2}\Psi_\lambda(c^\star(\lambda))\big)
\ge\mu>0$. Then $c^\star(\cdot)$ is a $C^{1}$, locally unique branch on
$[\lambda_{\min},\lambda_{\mathrm{cvx}}]$, and $c^\star(\lambda_{\min})$
is a strict local minimizer of $\Psi_{\lambda_{\min}}$ continuously
connected to the certified convex global. It is \emph{the} global
minimizer iff no minimizer on a disconnected branch has lower value ---
a condition the run can \emph{test} (multistart / monitoring) but not
\emph{prove} a priori.
\end{theorem}

\begin{proof}
$\partial_c\Theta=\nabla^{2}\Psi_\lambda$ is nonsingular along the path
by hypothesis, so the implicit function theorem
\cite{ortega1970,allgower1990} yields a unique $C^{1}$ branch
$c^\star(\lambda)$ on the interval, with no fold (a fold requires
$\sigma_{\min}\to0$). At $\lambda_{\mathrm{cvx}}$ the branch coincides
with the unique global of Theorem~\ref{thm:convex}; by continuity it
remains a strict local minimizer (positive-definite Hessian) throughout.
Global optimality at $\lambda_{\min}$ would require excluding a distinct
branch of lower value; nonconvexity permits such a branch, and its
absence is exactly the part not certifiable from
\eqref{eq:KT}--\eqref{eq:lamcvx} (Example~\ref{ex:scalar}).
\end{proof}

\begin{remark}[Honest scope]
Theorem~\ref{thm:path} is a genuine \emph{per-run, a-posteriori}
certificate: monitor $\sigma_{\min}$ along the homotopy; if it stays
$\ge\mu$, the produced solution is certified \emph{connected to the
certified convex global}; if a fold is detected, the method reports
``uncertified'' rather than a silent wrong answer. This is the strongest
guarantee available without assuming TCC or paying exponential cost
(\S\ref{sec:ceiling}).
\end{remark}

\section{Identifiability and the impossibility wall}\label{sec:ident}

Let $M(c)=J(c)^{\trc}J(c)$ be the Fisher information \cite{rao1945}.
Curvature (TCC, $K_T$) is irrelevant if $M$ is rank deficient: that is a
\emph{rank}, not a \emph{shape}, pathology, and it is fatal.

\begin{theorem}[Cram\'er--Rao impossibility]\label{thm:cr}
Let $0\neq v\in\ker F'(c^\star)\cap E$, $E$ the EA-coefficient subspace.
Then: \textup{(i)} $M(c^\star)$ is singular with $M(c^\star)v=0$, and the
Gauss--Newton (information) model of the misfit is flat along $v$, so the
parameter combination $\langle v,\cdot\rangle$ is not locally identifiable
from the data; \textup{(ii)} no unbiased estimator has finite variance
along $v$. No choice of $\lambda$, seed, or solver creates the missing
information.
\end{theorem}

\begin{proof}
(i) $M(c^\star)=F'(c^\star)^{\trc}F'(c^\star)$, so
$M(c^\star)v=F'(c^\star)^{\trc}\big(F'(c^\star)v\big)=0$. The
Gauss--Newton model $c\mapsto\tfrac12\norm{F(c^\star)+F'(c^\star)(c-
c^\star)}^{2}$ is constant along $c^\star+\spn\{v\}$ (since
$F'(c^\star)v=0$); equivalently the data carry no first-order information
about $\langle v,\cdot\rangle$. (ii) The Cram\'er--Rao inequality states
that any unbiased estimator obeys
$\operatorname{Cov}(\hat c)\succeq M^{+}$ on $\operatorname{range}M$ and
is \emph{unconstrained} (infinite variance) on $\ker M$
\cite{rao1945,walter1997}; with $Mv=0$,
$v^{\trc}\operatorname{Cov}(\hat c)v$ is unbounded. A Tikhonov term
biases the estimate toward the prior along $v$ but injects no data
information into $\ker M$, so it cannot restore identifiability --- only a
sensor change can (\S\ref{sec:multi}).
\end{proof}

\begin{remark}
TCC failure ($\eta\ge1$, curvature) and identifiability failure
($\ker M\neq\{0\}$, rank) are orthogonal. The former is repaired by
continuation/reformulation; the latter is irreparable for that sensor
configuration --- it must be reported, not optimized around.
\end{remark}

\section{Multisensor fusion}\label{sec:multi}

Stack heterogeneous channels $F=(F_1;\dots;F_S)$ (radar:
$[\rho,\mathrm{az},\mathrm{el},\dot\rho]$, nonlinear; optics: angles,
nonlinear; position fix: linear in $r$; accelerometer: affine in $a$).
Two effects are decisive and provable.

\begin{lemma}[A near-linear channel can only lower $\eta$]\label{lem:stack}
Let $F=(F_1;F_2)$ with $F_2$ affine. If $F_1$ satisfies the TCC with
constant $\eta_1$ on $B_\rho$, then the stacked $F$ satisfies it with
\[
\eta\;\le\;\eta_1\,\frac{\norm{\Delta F_1}}
{\sqrt{\norm{\Delta F_1}^{2}+\norm{\Delta F_2}^{2}}}\;\le\;\eta_1,
\]
where $\Delta F_j=F_j(c)-F_j(\tilde c)$; strict decrease whenever
$\Delta F_2\neq0$.
\end{lemma}

\begin{proof}
$F_2$ affine $\Rightarrow F_2(c)-F_2(\tilde c)-F_2'(c-\tilde c)=0$, so the
linearization remainder of the stack equals that of $F_1$:
$\norm{R(F)}=\norm{R(F_1)}\le\eta_1\norm{\Delta F_1}$. The denominator in
\eqref{eq:tcc} is
$\norm{\Delta F}=\sqrt{\norm{\Delta F_1}^{2}+\norm{\Delta F_2}^{2}}
\ge\norm{\Delta F_1}$. Divide.
\end{proof}

Thus fusing a position fix or an accelerometer --- channels affine (or
nearly so) in the unknowns --- \emph{provably} contracts the cone
constant and can move a TCC-failing radar-only problem into the
TCC-holding regime where Theorem~\ref{thm:irgn} and the certificates of
\S\S\ref{sec:convex}--\ref{sec:path} apply.

\begin{proposition}[Two radars restore identifiability]\label{prop:tworad}
In the degenerate single-radar, range-rate-on configuration, the
instantaneous measurement directions are the line-of-sight unit vector
and its derivative; the EA component the arc geometry leaves orthogonal
to the excited \textup{(LOS, $\dot{\text{LOS}}$)} directions lies in
$\ker F'\cap E$ (Theorem~\ref{thm:cr} then applies). Adding a second
radar with non-collinear lines of sight makes the per-epoch measurement
directions span $\R^{3}$, so generically $\ker F'\cap E=\{0\}$:
identifiability is restored by parallax.
\end{proposition}
\begin{proof}[Sketch]
The radar Jacobian rows are the LOS unit vector and its derivative; their
span has dimension $\le2$ per epoch for one station, $=3$ for two
non-collinear stations. The persistent EA subspace excited over the arc
then meets $\ker F'$ trivially; a standard observability/rank argument
\cite{walter1997} completes it.
\end{proof}

\begin{remark}[The $\sigma$-ladder is a TCC-restoring continuation]
Down-weighting the curved, ill-conditioned range-rate channel and
reintroducing it as the iterate approaches truth is exactly a homotopy
that keeps the \emph{effective} stacked $\eta<1$ along the path
(Lemma~\ref{lem:stack} with a schedule), i.e.\ an instance of
\S\ref{sec:path} with continuation in the data weighting rather than in
$\lambda$. Its path-dependence is the requirement to stay on the regular
branch of Theorem~\ref{thm:path}.
\end{remark}

\section{Reformulations that restore the hypotheses}\label{sec:reform}

Rather than \emph{assume} TCC, one can change $F$ so the certifiable
conditions \emph{hold}.

\begin{lemma}[Multiple shooting eliminates the $e^{L_gT}$ amplification]
\label{lem:ms}
Partition $[0,T]$ into $S$ segments of length $\Delta=T/S$ and optimize
jointly over per-node states and EA coefficients with continuity
constraints. The per-segment forward map satisfies
Lemma~\ref{lem:KT} with $T$ replaced by $\Delta$, hence per-segment
$K_\Delta=Ce^{L_g\Delta}\to C$ as $S\to\infty$. The exponential
nonlinearity amplification is removed; the coupling is by \emph{linear}
constraints.
\end{lemma}
\begin{proof}
Apply Lemma~\ref{lem:KT} on each segment $[t_{s},t_{s+1}]$ with arc
length $\Delta$; the Gr\"onwall factor is $e^{L_g\Delta}$. Continuity
$x(t_{s+1}^-)=x(t_{s+1}^+)$ is affine in the node states.
\end{proof}

This does not convexify the global problem (the constraint manifold is
nonconvex), but it deletes the dominant source of nonconvexity
(Example~\ref{ex:scalar}'s $\omega\ell$ with $\ell$ now $\Delta$), so the
SQP subproblems are well-conditioned and the certificate of
\S\ref{sec:convex} applies per segment. Pseudo-position pre-processing
(triangulating radar to an ECEF point per epoch) makes $\mathcal H$
affine, invoking Lemma~\ref{lem:stack}: it provably lowers $\eta$.

\section{The decision procedure}\label{sec:proc}

The preceding theorems compose into a constructive procedure whose every
branch is proved.

\begin{enumerate}
\item \textbf{Identifiability.} Estimate $\sigma_{\min}$ of the EA block
of $M$ on $\mathcal C$. If $\approx0$: \emph{stop} --- the global is a
manifold, recovery impossible for this configuration
(Theorem~\ref{thm:cr}). Fix by adding a non-collinear sensor
(Proposition~\ref{prop:tworad}) or removing the degenerate channel.
\item \textbf{Curvature.} From the force/state bounds compute $K_T$
(Lemma~\ref{lem:KT}) and $\lambda_{\mathrm{cvx}}$ \eqref{eq:lamcvx}.
\item \textbf{Certified surrogate.} Solve at
$\lambda\ge\lambda_{\mathrm{cvx}}$ from any seed; verify
$\nabla^{2}\Psi_\lambda\succ0$. This solve is certified global for the
surrogate (Theorem~\ref{thm:convex}, Corollary~\ref{cor:surr}).
\item \textbf{Continuation.} Decrease $\lambda\to\lambda_{\min}$,
monitoring $\sigma_{\min}(\nabla^{2}\Psi_\lambda)\ge\mu$. If maintained:
the result is certified connected to the surrogate global
(Theorem~\ref{thm:path}). If a fold is detected: report
\emph{uncertified}.
\item \textbf{Reformulate on failure.} Multiple shooting / collocation
(Lemma~\ref{lem:ms}) or pseudo-position (Lemma~\ref{lem:stack}) restores
the hypotheses; re-enter at step~2. Selection among admissible runs uses
a cross-validation criterion consistent under serial dependence
\cite{racine2000,burman1994,arlot2010}.
\end{enumerate}

\noindent What is proved: steps 1, 3, 5 carry full a-priori proofs; step
4 carries a per-run a-posteriori proof; the only unprovable element is
the \emph{absence of a disconnected lower branch}, which is \emph{tested}
(multistart agreement) and, by \S\ref{sec:ceiling}, cannot be eliminated
in general.

\section{The ceiling: what cannot be guaranteed, and why}
\label{sec:ceiling}

\begin{theorem}[No free unconditional global certificate]\label{thm:ceil}
For the operational problem \eqref{eq:Psi} at $\lambda_{\min}$ there is
no algorithm returning a certified global minimizer in time polynomial in
$p$ unless $\mathrm{P}=\mathrm{NP}$.
\end{theorem}
\begin{proof}[Sketch]
General nonconvex quadratic / nonlinear least-squares global optimization
is NP-hard \cite{murty1987}; \eqref{eq:Psi} contains such instances as
the linearization is varied. Assumption-free certified methods exist ---
interval/Lipschitz branch-and-bound and the moment--sum-of-squares
hierarchy \cite{neumaier2004,lasserre2001} --- but their cost is
exponential in $p$ ($p\sim10^{2}$--$10^{3}$ here), and validated ODE
enclosure suffers interval growth $\propto e^{L_gT}$ --- the \emph{same}
quantity (Example~\ref{ex:scalar}) that bounds $K_T$ and breaks TCC.
\end{proof}

Every route --- TCC's $\eta$, the {\L}ojasiewicz $\mu$, the contraction
threshold $\lambda_{\mathrm{cvx}}$, the interval-wrapping rate, the
path-fold likelihood --- is governed by the single intrinsic quantity
$\sim e^{L_gT}$. One therefore cannot escape it by changing proof
technique; one can only (i) shrink it by reformulation
(\S\ref{sec:reform}, Lemma~\ref{lem:ms}), (ii) certify a regularized
surrogate unconditionally (\S\ref{sec:convex}), (iii) certify a single
run a posteriori (\S\ref{sec:path}), or (iv) report impossibility when
the obstruction is rank, not curvature (\S\ref{sec:ident}). That
quadruple \emph{is} the answer to ``how do I ensure, and prove, the
global EA?'': a conditional, constructive, per-run certificate, with a
sharp and honest statement of its limits.

\begin{thebibliography}{99}
\bibitem{tikhonov1977} A.~N.~Tikhonov, V.~Y.~Arsenin.
 \textit{Solutions of Ill-Posed Problems}. Winston, 1977.
\bibitem{engl1996} H.~W.~Engl, M.~Hanke, A.~Neubauer.
 \textit{Regularization of Inverse Problems}. Kluwer, 1996.
\bibitem{hairer1993} E.~Hairer, S.~P.~N\o rsett, G.~Wanner.
 \textit{Solving ODEs I}, 2nd ed. Springer, 1993.
\bibitem{ortega1970} J.~M.~Ortega, W.~C.~Rheinboldt. \textit{Iterative
 Solution of Nonlinear Equations in Several Variables}. Academic, 1970.
\bibitem{deuflhard2004} P.~Deuflhard. \textit{Newton Methods for
 Nonlinear Problems}. Springer, 2004.
\bibitem{nocedal2006} J.~Nocedal, S.~J.~Wright. \textit{Numerical
 Optimization}, 2nd ed. Springer, 2006.
\bibitem{scherzer1995} O.~Scherzer. Convergence criteria of iterative
 methods\dots\ \textit{J.~Math.\ Anal.\ Appl.} 194:911--933, 1995.
\bibitem{hanke1995} M.~Hanke, A.~Neubauer, O.~Scherzer. A convergence
 analysis of the Landweber iteration for nonlinear ill-posed problems.
 \textit{Numer.\ Math.} 72:21--37, 1995.
\bibitem{kaltenbacher2008} B.~Kaltenbacher, A.~Neubauer, O.~Scherzer.
 \textit{Iterative Regularization Methods for Nonlinear Ill-Posed
 Problems}. de~Gruyter, 2008.
\bibitem{bakushinskii1992} A.~B.~Bakushinskii. The problem of the
 convergence of the iteratively regularized Gauss--Newton method.
 \textit{Comput.\ Math.\ Math.\ Phys.} 32:1353--1359, 1992.
\bibitem{polyak1963} B.~T.~Polyak. Gradient methods for minimizing
 functionals. \textit{Zh.\ Vychisl.\ Mat.} 3:643--653, 1963.
\bibitem{lojasiewicz1963} S.~{\L}ojasiewicz. Une propri\'et\'e
 topologique des sous-ensembles analytiques r\'eels. 1963.
\bibitem{karimi2016} H.~Karimi, J.~Nutini, M.~Schmidt. Linear convergence
 of gradient and proximal-gradient methods under the
 Polyak--{\L}ojasiewicz condition. \textit{ECML PKDD}, 2016.
\bibitem{nesterov2004} Y.~Nesterov. \textit{Introductory Lectures on
 Convex Optimization}. Kluwer, 2004.
\bibitem{rao1945} C.~R.~Rao. Information and accuracy attainable in the
 estimation of statistical parameters. \textit{Bull.\ Calcutta Math.\
 Soc.} 37:81--89, 1945.
\bibitem{walter1997} E.~Walter, L.~Pronzato. \textit{Identification of
 Parametric Models from Experimental Data}. Springer, 1997.
\bibitem{allgower1990} E.~L.~Allgower, K.~Georg. \textit{Numerical
 Continuation Methods}. Springer, 1990.
\bibitem{racine2000} J.~Racine. Consistent cross-validatory
 model-selection for dependent data: hv-block cross-validation.
 \textit{J.~Econometrics} 99:39--61, 2000.
\bibitem{burman1994} P.~Burman, E.~Chow, D.~Nolan. A cross-validatory
 method for dependent data. \textit{Biometrika} 81:351--358, 1994.
\bibitem{arlot2010} S.~Arlot, A.~Celisse. A survey of cross-validation
 procedures for model selection. \textit{Statist.\ Surv.} 4:40--79, 2010.
\bibitem{murty1987} K.~G.~Murty, S.~N.~Kabadi. Some NP-complete problems
 in quadratic and nonlinear programming. \textit{Math.\ Program.}
 39:117--129, 1987.
\bibitem{neumaier2004} A.~Neumaier. Complete search in continuous global
 optimization and constraint satisfaction. \textit{Acta Numerica}
 13:271--369, 2004.
\bibitem{lasserre2001} J.~B.~Lasserre. Global optimization with
 polynomials and the problem of moments. \textit{SIAM J.~Optim.}
 11:796--817, 2001.
\end{thebibliography}

\end{document}
